% Standalone chunk: Corollary on robustness to approximate low rank
% + variance-misspecification remark (POST-FIX, Round 2).
% Source: conf.tex lines ~563-651 (post-fix).
% Shared preamble for standalone chunks. Do not compile directly; \input{} it.
\documentclass{article}
\usepackage[utf8]{inputenc}
\usepackage[T1]{fontenc}
\usepackage{lmodern}
\usepackage[margin=1in]{geometry}
\usepackage{amsmath,amssymb,amsfonts,amsthm,bm,mathtools}
\usepackage{enumitem}
\usepackage{microtype}
\usepackage{xcolor}
\usepackage{natbib}
\usepackage{booktabs}
\usepackage{hyperref}

\newcommand{\R}{\mathbb{R}}
\newcommand{\N}{\mathbb{N}}
\newcommand{\E}{\mathbb{E}}
\newcommand{\Prob}{\mathbb{P}}
\newcommand{\op}{\mathrm{op}}
\DeclareMathOperator{\sign}{sign}
\newcommand{\cA}{\mathcal{A}}
\newcommand{\cF}{\mathcal{F}}
\newcommand{\cH}{\mathcal{H}}
\newcommand{\cT}{\mathcal{T}}
\newcommand{\cI}{\mathcal{I}}
\newcommand{\cW}{\mathcal{W}}
\newcommand{\cTprobe}{\mathcal{T}_{\mathrm{probe}}}
\newcommand{\cE}{\mathcal{E}}
\newcommand{\norm}[1]{\left\lVert #1 \right\rVert}
\newcommand{\tilO}{\widetilde{\mathcal{O}}}
\newcommand{\DynReg}{\mathrm{DynReg}}
\newcommand{\rank}{\mathrm{rank}}
\newcommand{\range}{\mathrm{range}}
\DeclareMathOperator{\TopEig}{TopEig}
\DeclareMathOperator*{\argmax}{arg\,max}
\DeclareMathOperator{\tr}{tr}

\newtheorem{assumption}{Assumption}
\newtheorem{theorem}{Theorem}
\newtheorem{corollary}{Corollary}
\newtheorem{proposition}{Proposition}
\newtheorem{lemma}{Lemma}
\newtheorem{remark}{Remark}


\title{Chunk: Cor.\ misspec + variance-misspec remark\\(post-fix, Round 2 verification)}
\author{(excerpt for standalone review)}
\date{}

\begin{document}
\maketitle

\textbf{Setup context.}
\begin{itemize}[leftmargin=*,itemsep=1pt,topsep=1pt]
\item SPSC's probe-moment concentration: Cor.\ projector\_conf gives
$\varepsilon_k\le C_{\mathrm{sub}}\sqrt{\log(2d/\delta)/m_k}+\Delta_\sigma$
with $\Delta_\sigma=4(|\delta_\sigma|L^2+2L^3\epsilon_\times)/\lambda_{\min}$
(conservative, any probe distribution).
\item Main upper bound (Thm.\ 1):
$\DynReg_T^{(c)}=\tilO(r\sqrt{KT})+\tilO(K^{1/3}T^{2/3})+O(WV)+O(T\Delta_\sigma)$.
\item Exploit-regret decomposition has a subspace-mismatch term
$R_\cA S_w\varepsilon_k n_k$ per segment (equal-rank-projector identity
$\|(I-\widehat P)B^\star\|_\op\le\varepsilon_k$).
\end{itemize}

\textbf{What changed since Round~1.} No substantive revisions to this
chunk. Minor wording clarifications around the $\delta_\sigma$-bias
shift (this remark is stable across R1 and R2 because the sharpened
bias derivation is housed in chunk~A.2's Lem.~G\_unbiased\_conf and
Rem.~bias\_floor\_sharp; this chunk just points to them).

\section*{Corollary (approximate low rank)}

\begin{corollary}[Robustness to approximate low rank]
\label{cor:misspec}
Suppose $\theta_t = B_k^\star w_t + \theta_t^{\perp}$ inside $\cI_k$,
with $B_k^\star$ orthonormal as before and
$\sup_{t\in\cI_k}\|\theta_t^\perp\|_2\le\epsilon_k^\perp \le S_w$,
and assume the strengthened eigengap condition
$\lambda_{\min}\ge 4\cdot C_{\mathrm{sub}}\sqrt{\log(2d/\delta)/m_k}
+ 12\,S_w\,\epsilon_k^\perp$.
Then the bound of Thm.~1 holds with the additive term
\[
O\!\Bigl(R_\cA\sum_k \epsilon_k^\perp\,\bigl(1 + S_w^2/\lambda_{\min}\bigr)\, n_k\Bigr)
\]
added to the right-hand side. When $\epsilon_k^\perp \le c\,\lambda_{\min}/S_w^2$,
the additive term is $O(R_\cA\sum_k \epsilon_k^\perp n_k)$ and the
approximate-rank penalty is linear in $\epsilon^\perp$ per round.
\end{corollary}
\begin{proof}[Sketch]
Decompose the predictable second moment using $\theta_t=B_k^\star w_t+\theta_t^\perp$:
$\widetilde M_t = B_k^\star\,\mathbb E[w_tw_t^\top\mid\cH_{t-1}]\,(B_k^\star)^\top + E_t$,
where the perturbation
$E_t = B_k^\star\,\mathbb E[w_t(\theta_t^\perp)^\top\mid\cH_{t-1}]
+ \mathbb E[\theta_t^\perp w_t^\top\mid\cH_{t-1}]\,(B_k^\star)^\top
+ \mathbb E[\theta_t^\perp(\theta_t^\perp)^\top\mid\cH_{t-1}]$
satisfies $\|E_t\|_\op\le 2 S_w\,\epsilon_k^\perp + (\epsilon_k^\perp)^2
\le 3 S_w\,\epsilon_k^\perp$ (using $\epsilon_k^\perp\le S_w$).

Davis--Kahan applied to $\widehat M_k$ under the strengthened eigengap
condition gives
$\varepsilon_k\le 4 C_{\mathrm{sub}}\sqrt{\log(2d/\delta)/m_k}/\lambda_{\min}
+ 12\,S_w\,\epsilon_k^\perp/\lambda_{\min}$.

The regret decomposition contributes
$R_\cA(I-\widehat P_{t-1})\theta_t = R_\cA(I-\widehat P_{t-1})
(B_k^\star w_t + \theta_t^\perp)$, whose norm is bounded by
$R_\cA(\varepsilon_k S_w + \epsilon_k^\perp)$ per round (using
$\|(I-\widehat P)B_k^\star\|_\op\le\|P_k^\star-\widehat P\|_\op$). Summing
over $E_k$:
$R_\cA n_k (\varepsilon_k S_w + \epsilon_k^\perp)
\le R_\cA n_k \epsilon_k^\perp(12 S_w^2/\lambda_{\min} + 1) + (\text{statistical piece})$.
The statistical piece is the same one optimized in Thm.~1 and
contributes to the leading $K^{1/3}T^{2/3}$ term.
\end{proof}

\section*{Remark (variance misspecification)}

\begin{remark}[Variance misspecification: two bounds]
\label{rem:variance_misspec}
Write $\delta_\sigma=\hat\sigma^2-\sigma_\varepsilon^2$. The bias floor
in Thm.~1 includes the term $|\delta_\sigma|L^2/\lambda_{\min}$
because the general Davis--Kahan chain bounds the operator-norm bias
of $\widehat M_k$ conservatively, without exploiting the structure of
that bias.

For \emph{isotropic Gaussian probes}, however, the $\delta_\sigma$-component
of the bias is
$-\delta_\sigma\,\mathcal K^{-1}(\mathbb E[uu^\top])=-\delta_\sigma/(d+2)\cdot I_d$
(by Lem.~K\_inverse)---a \emph{scaled identity}---which shifts the
spectrum of $\widehat M_k$ uniformly but does not rotate eigenvectors.
The top-$r$ subspace of $\widehat M_k$ is therefore unchanged, so the
bias propagation through Davis--Kahan (which only responds to
eigenvector-rotating perturbations) is actually driven only by the
$L^3\epsilon_\times$ part of the bias, not $|\delta_\sigma|L^2$.

A sharper statement of Thm.~1 therefore has
$\Delta_\sigma^{\mathrm{sharp}}:=8 L^3\epsilon_\times/\lambda_{\min}$
in place of $\Delta_\sigma$. The conservative $\Delta_\sigma$ in the
theorem statement is valid for any probe distribution; the sharper
$\Delta_\sigma^{\mathrm{sharp}}$ is specific to isotropic Gaussian probes.
We always have $\Delta_\sigma^{\mathrm{sharp}}\le\Delta_\sigma$.
\end{remark}

\end{document}
